\documentclass[11pt]{article}
\usepackage{amsmath,amssymb,booktabs,geometry,graphicx,hyperref,url}
\usepackage{subcaption,placeins}
\hypersetup{hidelinks}
\geometry{margin=1in}
\setlength{\parskip}{0.45em}
\setlength{\parindent}{0pt}
\graphicspath{{./}{../}}
\title{Stage 6: Surrogate Dataset Construction}
\author{Free-Boundary-Aware Risk Surfaces for American Options}
\date{\today}
\begin{document}
\maketitle

\section{Stage 6: Surrogate Dataset Construction}

\subsection{Purpose}
This stage constructs a supervised dataset for later surrogate experiments
using audited CN/PSOR labels. It is a data-engineering and governance stage:
the central question is which solver regimes are sufficiently reliable to
become labels. No neural network is trained, selected, or evaluated here.

\subsection{Method}
The dataset consumes the Stage 5 application outputs:
\[
\texttt{risk\_metrics.csv},\quad
\texttt{regime\_flags.csv},\quad
\texttt{boundary\_summary.csv}.
\]
Only American put and dividend-paying American call regimes with no diagnostic failure and no delicate-regime flag are included as clean labels. No-dividend American call controls and all flagged regimes are retained in \texttt{excluded\_regimes.csv} with explicit reasons.

The primary split is regime-level rather than pointwise. With seed 7, the 223
clean regimes are split into 156 training, 33 validation, and 34 test regimes.
All spot-time nodes from one parameter regime remain in the same partition.
This prevents nearly identical neighboring nodes from appearing on both sides
of the evaluation boundary and creating an unrealistically easy test.

\subsection{Mathematical Formulation}
For each clean regime, the finite-difference solver produces a surface $U(S,\tau)$, payoff $\Phi(S)$, continuation premium
\[
C(S,\tau)=\max(U(S,\tau)-\Phi(S),0),
\]
finite-difference Delta, and finite-difference Gamma. Each node contributes the feature vector
\[
\left(\log(S/K),\;\tau/T,\;r,\;q,\;\sigma,\;T,\;\mathbf{1}_{\mathrm{call}}\right).
\]
The label vector is
\[
\left(U/K,\;C/K,\;\Phi/K,\;\Delta,\;K\Gamma\right).
\]
Greek labels are accompanied by strict masks and exclusion masks for the payoff kink, free-boundary band, and near-zero spot region.

\subsection{Implementation Details}
The main workflow is \texttt{experiments/05\_surrogate\_dataset.py}. Dataset sample construction is implemented in \texttt{src/data/dataset\_builder.py}. Regime-level train/validation/test splitting is implemented in \texttt{src/data/splits.py}.

The generated dataset is \path{data/processed/surrogate_dataset.npz}. Its arrays
are grouped as follows:
\begin{itemize}
  \item core labels: \texttt{features}, \texttt{labels}, \texttt{boundary},
    and \texttt{exercise\_indicator};
  \item diagnostic masks: \texttt{strict\_greek\_mask}, \texttt{near\_kink},
    \texttt{near\_boundary}, and \texttt{near\_zero\_spot};
  \item partition metadata: \texttt{regime\_key} and \texttt{split};
  \item audit metadata: \texttt{diagnostic\_status}, \texttt{source\_stage},
    and \texttt{solver\_config\_hash}.
\end{itemize}

\subsection{Validation and Diagnostics}
\begin{table}[h]
\centering
\begin{tabular}{lr}
\toprule
Quantity & Value\\
\midrule
Clean regimes & 223\\
Excluded regimes & 65\\
Total samples & 3,264,943\\
Feature dimension & 7\\
Label dimension & 5\\
Train regimes & 156\\
Validation regimes & 33\\
Test regimes & 34\\
\bottomrule
\end{tabular}
\caption{Clean-label dataset construction summary after Stage 5 diagnostic filtering. Counts are reported by regime and sample so that the much larger pointwise sample total is not confused with independent market scenarios.}
\end{table}

The excluded regimes consist of 33 control-only regimes, 29 diagnostic-failure/delicate regimes, and 3 no-dividend control regimes that also carried diagnostic flags. No regime key appears in more than one split.

\begin{table}[h]
\centering
\begin{tabular}{lrr}
\toprule
Mask & Samples & Fraction\\
\midrule
Exercise indicator & 1,766,509 & 0.5411\\
Strict Greek mask & 3,157,110 & 0.9670\\
Near kink & 26,983 & 0.0083\\
Near boundary & 26,983 & 0.0083\\
Near zero spot & 53,966 & 0.0165\\
\bottomrule
\end{tabular}
\caption{Coverage of exercise and Greek-quality masks across the 3,264,943 clean samples. The strict Greek mask excludes the payoff kink, free-boundary band, and near-zero-spot region.}
\end{table}

\subsection{Results}
The dataset contains normalized price, nonnegative continuation premium, payoff, Delta, and $K\Gamma$ labels. The continuation premium ranges from 0 to about 0.2634. The strict Greek mask covers about 96.70\% of nodes, leaving explicit exclusions for delicate Greek zones.

\begin{figure}[h]
\centering
\includegraphics[width=0.45\linewidth]{results/figures/stage06_surrogate_dataset/price_over_K_hist.png}
\includegraphics[width=0.45\linewidth]{results/figures/stage06_surrogate_dataset/premium_over_K_hist.png}
\caption{Distributions of normalized price \(U/K\) and continuation premium \(C/K\) in the clean dataset. These plots describe label coverage and do not constitute surrogate-performance evidence.}
\end{figure}

\begin{figure}[h]
\centering
\includegraphics[width=0.45\linewidth]{results/figures/stage06_surrogate_dataset/delta_hist.png}
\includegraphics[width=0.45\linewidth]{results/figures/stage06_surrogate_dataset/K_gamma_hist.png}
\caption{Distributions of finite-difference Delta and scaled Gamma \(K\Gamma\). Gamma remains a masked diagnostic label because second derivatives are least stable near kinks and free boundaries.}
\end{figure}

% Auto-generated by experiments/19_report_visualization_enrichment.py
\section*{Visual Evidence Gallery}
\begin{figure}[!htbp]
\centering
\includegraphics[width=0.84\linewidth]{results/figures/report_enrichment/stage06_surrogate_dataset_dataset_flow.pdf}
\caption{Clean supervised labels are constructed only after CN/PSOR validation, application diagnostics, and explicit exclusion of control or flagged regimes.}
\end{figure}

\begin{figure}[!htbp]
\centering
\includegraphics[width=0.70\linewidth]{results/figures/report_enrichment/stage06_surrogate_dataset_regime_split.pdf}
\caption{Training, validation, and test partitions are assigned by regime key so neighboring nodes from one PSOR surface cannot leak across splits.}
\end{figure}

\begin{figure}[!htbp]
\centering
\begin{subfigure}[t]{0.48\linewidth}
\centering
\includegraphics[width=\linewidth]{results/figures/report_enrichment/stage06_surrogate_dataset_label_summary.pdf}
\caption{Label summaries document U over K, continuation premium C over K, Delta, and scaled Gamma K Gamma before neural training; Gamma remains mask-qualified.}
\end{subfigure}
\hfill
\begin{subfigure}[t]{0.48\linewidth}
\centering
\includegraphics[width=\linewidth]{results/figures/report_enrichment/stage06_surrogate_dataset_mask_summary.pdf}
\caption{Mask composition records how much data remains for strict Greek evaluation.}
\end{subfigure}
\end{figure}

\FloatBarrier


\subsection{Failure Modes and Limitations}
This dataset excludes the 32 Stage 5 diagnostic-failure regimes rather than
repairing them and excludes no-dividend call controls from the clean ML
dataset. Because three controls also carry diagnostic flags, the unique
exclusion table contains 65 regimes rather than the sum of disjoint categories.
Later surrogate experiments must therefore state that their training envelope
is the clean, diagnostic-passing American put and dividend-paying call
envelope. Gamma labels are included only with masks; their presence does not
imply that unmasked Gamma is stable or that a Gamma surrogate is validated.

\subsection{Conclusion and Next Step}
The Stage 6 gate passes for beginning baseline surrogate experiments on the clean dataset. It does not provide surrogate-performance evidence. The next stage should train baseline models and audit their obstacle, monotonicity, convexity, boundary, and Greek failures without hiding weak regimes.

\end{document}
